% GLUE THAT SHRINKS WITHOUT LIMIT IN A PARAGRAPH.
%
% It would let a paragraph of any length fit on one line, so the
% reference says so and takes the shrinkage as finite:
%
%   ! Infinite glue shrinkage found in a paragraph.
%   The paragraph just ended includes some glue that has
%   infinite shrinkability, e.g., `\hskip 0pt minus 1fil'.
%   Such glue doesn't belong there---it allows a paragraph
%   of any length to fit on one line. But it's safe to proceed,
%   since the offensive shrinkability has been made finite.
%
% THE AMOUNT IS KEPT AND THE ORDER DROPPED: `\hskip 3pt minus 1fil'
% comes out of the line as `\glue 3.0 minus 1.0'. A \rightskip treated
% that way stays treated -- `\showthe\rightskip' afterwards reads
% `0.0pt minus 1.0pt' -- and it is not an assignment, so no save is made
% and no group puts the fil back.
%
% ONCE A PARAGRAPH, however many offenders there are and whether they
% were skip parameters or glue in the list.
%
% WHERE THE LINE COMES depends on which it was, and \tracingparagraphs
% shows it:
%
%   \rightskip     the fault, then `@firstpass'
%   glue in the list   `@firstpass', the breakpoint before the glue,
%                      THEN the fault
%
% -- because the skips are looked at before the passes start and the
% list's own glue is met during the first pass, at the glue. trip line
% 210 is the first kind.
%
% AND THE TRACE IS CLOSED BEFORE THE FAULT AND OPENED AGAIN AFTER, which
% leaves ONE MORE BLANK LINE above the `!' than there would be with
% \tracingparagraphs off. That holds for both kinds -- the fault before
% the first pass gets it too, because the trace's diagnostic is already
% open by then, which is also why `@firstpass' comes out AFTER that
% fault rather than before it.
\tracingonline=1 \hsize=100pt \parindent=0pt \showboxdepth=5 \showboxbreadth=10
\message{[A finite shrink draws nothing]}
\noindent A\hskip 0pt minus 5pt B\par
\message{[B glue in the list]}
\setbox0=\vbox{\noindent A\hskip 3pt minus 1fil B\par}
\showbox0
\message{[C two offenders, one report]}
\noindent A\hskip 0pt minus 1fil B\hskip 0pt minus 1fil C\par
\message{[D a skip parameter]}
\rightskip=0pt minus 1fil
\setbox1=\vbox{\noindent A B\par}
\showbox1
\showthe\rightskip
\rightskip=0pt
\message{[E where the line comes]}
\tracingparagraphs=1 \pretolerance=10000
\noindent A\hskip 3pt minus 1fil B\par
\rightskip=0pt minus 1fil
\noindent A B\par
\rightskip=0pt \tracingparagraphs=0
\message{[done]}
\end
